\chapter{Cholecystectomy}

\section*{Overview}
Cholecystectomy is surgical removal of the gallbladder, most often to treat symptomatic gallstones or inflammation of the gallbladder.

\section*{Why It Is Done}
It may be recommended for recurrent biliary colic, acute or chronic cholecystitis, gallstone pancreatitis, or selected gallbladder complications.

\section*{How It Is Performed}
Usually laparoscopically under general anesthesia, the surgeon clips and divides the cystic duct and artery and removes the gallbladder through small ports. Open surgery may be needed in complicated cases.

\section*{Risks and Recovery}
Bleeding, infection, bile leak, bile-duct injury, retained stones, blood clots, and anesthesia complications are possible. Most uncomplicated patients go home the same day or after an overnight stay and gradually resume activity.

\section*{References}
\begin{enumerate}
  \item MedlinePlus. Gallbladder Removal---Open. U.S. National Library of Medicine; 2025.
  \item Current Medical Diagnosis and Treatment. McGraw-Hill; 2025.
\end{enumerate}
\clearpage
